\chapter{Capsule Endoscopy}

\section*{Overview}
Capsule endoscopy is a procedure used to record images of the gastrointestinal tract for use in medicine. It utilizes a pill-sized capsule containing a wireless camera.

\section*{Alternative Names}
Pill cam, Wireless capsule endoscopy

\section*{Principle}
The patient swallows a small capsule containing a miniature camera, light source, battery, and transmitter. As peristalsis moves it through the digestive tract, it captures thousands of images that are transmitted to a recorder worn on the body, allowing visualization of the small intestine, which is hard to reach with standard endoscopes.
\section*{History}
Gavriel Iddan and Paul Swain co-invented the capsule endoscope in the late 1990s, and it was approved by the FDA in 2001.

\section*{Why It Is Done}
Ordered to find the source of bleeding in the small intestine (which cannot be reached by standard endoscopy), diagnose Crohn's disease, or detect small bowel tumors.

\section*{Is This a Primary or Secondary Test or Procedure}
Primary diagnostic tool for small bowel mucosal imaging and obscure GI bleeding.

\section*{How It Is Performed}
You swallow a vitamin-sized capsule containing a camera. As it travels naturally through your digestive tract, it transmits thousands of images to a recorder worn on your belt.

\section*{Preparation}
Fast for 12 hours before swallowing the capsule. A laxative bowel prep is usually required the night before.

\section*{Contraindications and Precautions}
Known or suspected gastrointestinal obstruction, stricture, or fistula. Severe swallowing disorders.

\section*{Risks}
Capsule retention (<2% of cases), which may require surgical or endoscopic removal.

\section*{Aftercare and Recovery}
You can eat a light snack after 4 hours. The capsule is disposable and passes naturally in your stool within 24 to 72 hours.

\section*{Results and Interpretation}
Images are downloaded from the recorder and reviewed by a gastroenterologist for small bowel ulcers, bleeding, or polyps.

\section*{Expected Results}
Normal small bowel mucosa; no active bleeding, ulcers, erosions, or vascular malformations.

\section*{Follow-up}
Double-balloon enteroscopy, CT enterography, or medical management.

\section*{References}
\begin{enumerate}
  \item MedlinePlus. Capsule Endoscopy. U.S. National Library of Medicine; 2025.
  \item Current Medical Diagnosis and Treatment. McGraw-Hill; 2025.
\end{enumerate}

\clearpage
